\section{Noetherian rings and finite constructions}
\label{am-ch7-noetherian}
\begin{definition}[Noetherian ring]
\label{am-def-noetherian-ring}
\source{atiyah-macdonald-commutative-algebra:page-0089}
A ring is Noetherian when its ideals satisfy the equivalent maximal condition,
ascending chain condition, or finite-generation condition.
\end{definition}
\begin{proposition}[Quotients]
\label{am-prop-7-1}
\dcref{7.1}
\source{atiyah-macdonald-commutative-algebra:page-0089}
Every homomorphic image of a Noetherian ring is Noetherian.
\uses{am-prop-6-6}
\end{proposition}
\begin{proof}
Any quotient is a quotient by the kernel, so apply Proposition~\ref{am-prop-6-6}.
\end{proof}
\begin{proposition}[Finite modules and finite algebras]
\label{am-prop-7-2}
\dcref{7.2}
\source{atiyah-macdonald-commutative-algebra:page-0089}
If $A$ is Noetherian and $B$ is finite as an $A$-module, then $B$ is a
Noetherian ring.
\end{proposition}
\begin{proof}
By Proposition~\ref{am-prop-6-5}, $B$ is Noetherian as an $A$-module, hence as
a $B$-module because every $B$-submodule is an $A$-submodule.
\uses{am-prop-6-5}
\end{proof}
\begin{proposition}[Localization]
\label{am-prop-7-3}
\dcref{7.3}
\source{atiyah-macdonald-commutative-algebra:page-0089}
Localization of a Noetherian ring at any multiplicative set is Noetherian.
\end{proposition}
\begin{proof}
Ideals of $S^{-1}A$ correspond order-preservingly to contracted ideals of $A$
by Proposition~\ref{am-prop-3-11}; the maximal condition transfers.
\uses{am-prop-3-11}
\end{proof}
\begin{corollary}[Localizations at primes]
\label{am-cor-7-4}
\dcref{7.4}
\source{atiyah-macdonald-commutative-algebra:page-0089}
If $A$ is Noetherian, $A_\p$ is Noetherian for every prime $\p$.
\uses{am-prop-7-3}
\end{corollary}
\begin{proof}
Take $S=A\setminus\p$.
\end{proof}
\begin{theorem}[Hilbert basis theorem]
\label{am-thm-7-5}
\dcref{7.5}
\source{atiyah-macdonald-commutative-algebra:page-0090}
If $A$ is Noetherian, then $A[x]$ is Noetherian.
\end{theorem}
\begin{proof}
For an ideal $\aideal\subseteq A[x]$, leading coefficients form a finitely
generated ideal of $A$.  Choose finitely many polynomials in $\aideal$ with
those leading coefficients and bounded degrees.  Subtract suitable multiples
to reduce every element of $\aideal$ to one of bounded degree.  The bounded
degree polynomials form a finite $A$-module, so their intersection with
$\aideal$ is finitely generated (Propositions~\ref{am-prop-6-5} and
\ref{am-prop-6-2}); together with the chosen leading polynomials these generate
$\aideal$.
\uses{am-prop-6-5,am-prop-6-2}
\end{proof}
\begin{corollary}[Polynomial rings]
\label{am-cor-7-6}
\dcref{7.6}
\source{atiyah-macdonald-commutative-algebra:page-0090}
$A[x_1,\ldots,x_n]$ is Noetherian when $A$ is.
\end{corollary}
\begin{proof}
Induct on $n$ using Theorem~\ref{am-thm-7-5}.
\uses{am-thm-7-5}
\end{proof}
\begin{corollary}[Finite-type algebras]
\label{am-cor-7-7}
\dcref{7.7}
\source{atiyah-macdonald-commutative-algebra:page-0090}
Every finite-type $A$-algebra over a Noetherian ring is Noetherian.
\uses{am-cor-7-6,am-prop-7-1}
\end{corollary}
\begin{proof}
It is a quotient of a polynomial ring, so use Corollary~\ref{am-cor-7-6} and
Proposition~\ref{am-prop-7-1}.
\end{proof}
\begin{proposition}[Artin--Tate lemma]
\label{am-prop-7-8}
\dcref{7.8}
\source{atiyah-macdonald-commutative-algebra:page-0090}
If $A\subseteq B\subseteq C$, $A$ is Noetherian, $C$ is finite type over $A$,
and $C$ is finite as a $B$-module (equivalently integral over $B$), then $B$ is
finite type over $A$.
\end{proposition}
\begin{proof}
Choose algebra generators $x_i$ of $C$ and module generators $y_j$.  The
coefficients expressing $x_i$ in the $y_j$ and products $y_iy_j$ generate an
$A$-algebra $B_0\subseteq B$ over which $C$ is finite.  Since $B_0$ is
Noetherian, its submodule $B\subseteq C$ is finite over $B_0$; adjoining module
generators of $B$ to the finite-type algebra $B_0$ generates $B$ as an
$A$-algebra.
\uses{am-prop-5-1,am-cor-5-2,am-cor-7-7,am-prop-6-2}
\end{proof}
\begin{theorem}[Weak Nullstellensatz]
\label{am-thm-7-9}
\dcref{7.9}
\source{atiyah-macdonald-commutative-algebra:page-0091}
If $k$ is a field and a finite-type $k$-algebra $E$ is a field, then $E$ is a
finite algebraic extension of $k$.
\end{theorem}
\begin{proof}
Choose a maximal algebraically independent subset among generators.  If it
were nonempty, write the remaining generators algebraic over the resulting
rational function field $F$.  Then $E$ is finite over $F$, so the Artin--Tate
lemma makes $F$ finite type over $k$.  A finite set of rational functions cannot
generate inverses of all irreducible polynomials in a positive-dimensional
polynomial ring, contradiction.  Thus all generators are algebraic and finite
generation makes the extension finite.
\uses{am-prop-7-8}
\end{proof}
\begin{corollary}[Residue fields]
\label{am-cor-7-10}
\dcref{7.10}
\source{atiyah-macdonald-commutative-algebra:page-0091}
For a finite-type $k$-algebra $A$ and maximal $\m$, $A/\m$ is finite algebraic
over $k$; if $k$ is algebraically closed, $A/\m\simeq k$.
\uses{am-thm-7-9}
\end{corollary}
\begin{proof}
Apply Theorem~\ref{am-thm-7-9} to $A/\m$.
\end{proof}

\section{Primary decomposition in Noetherian rings}
\label{am-ch7-primary}
\begin{definition}[Irreducible ideal]
\label{am-def-irreducible-ideal}
\source{atiyah-macdonald-commutative-algebra:page-0091}
An ideal is irreducible if $\aideal=\bideal\cap\cideal$ implies
$\aideal=\bideal$ or $\aideal=\cideal$.
\end{definition}
\begin{lemma}[Finite intersection of irreducibles]
\label{am-lem-7-11}
\dcref{7.11}
\source{atiyah-macdonald-commutative-algebra:page-0092}
Every ideal in a Noetherian ring is a finite intersection of irreducible ideals.
\end{lemma}
\begin{proof}
Otherwise the ideals failing this property have a maximal member $\aideal$.
It is reducible, $\aideal=\bideal\cap\cideal$ with both strictly larger; by
maximality each is a finite intersection of irreducibles, contradiction.
\end{proof}
\begin{lemma}[Irreducible ideals are primary]
\label{am-lem-7-12}
\dcref{7.12}
\source{atiyah-macdonald-commutative-algebra:page-0092}
Every irreducible ideal in a Noetherian ring is primary.
\end{lemma}
\begin{proof}
Pass to the quotient and suppose $(0)$ irreducible.  If $xy=0$ with $y\ne0$,
the ascending chain $\Ann(x)\subseteq\Ann(x^2)\subseteq\cdots$ stabilizes.
For a stable index $n$, $(x^n)\cap(y)=0$: if $a=by$ and $a=bx^n$, then
$bx^{n+1}=0$ implies $b\in\Ann(x^n)$ and hence $a=0$.  Irreducibility and
$y\ne0$ force $x^n=0$, so every zero-divisor is nilpotent.
\uses{am-prop-6-2}
\end{proof}
\begin{theorem}[Existence of primary decompositions]
\label{am-thm-7-13}
\dcref{7.13}
\source{atiyah-macdonald-commutative-algebra:page-0092}
Every ideal in a Noetherian ring has a primary decomposition.
\uses{am-lem-7-11,am-lem-7-12}
\end{theorem}
\begin{proof}
Combine Lemma~\ref{am-lem-7-11} with Lemma~\ref{am-lem-7-12}.
\end{proof}
\begin{proposition}[A power of the radical]
\label{am-prop-7-14}
\dcref{7.14}
\source{atiyah-macdonald-commutative-algebra:page-0092}
For every ideal $\aideal$ in a Noetherian ring, $\rad(\aideal)^n\subseteq\aideal$
for some $n$.
\end{proposition}
\begin{proof}
Choose generators $x_i$ of $\rad(\aideal)$ and exponents $x_i^{n_i}\in\aideal$.
For $n=\sum_i(n_i-1)+1$, every degree-$n$ monomial contains some
$x_i^{n_i}$, so lies in $\aideal$.
\end{proof}
\begin{corollary}[Nilradical nilpotent]
\label{am-cor-7-15}
\dcref{7.15}
\source{atiyah-macdonald-commutative-algebra:page-0092}
The nilradical of a Noetherian ring is nilpotent.
\uses{am-prop-7-14}
\end{corollary}
\begin{proof}
Take $\aideal=(0)$ in Proposition~\ref{am-prop-7-14}.
\end{proof}
\begin{corollary}[Characterizing primary ideals with maximal radical]
\label{am-cor-7-16}
\dcref{7.16}
\source{atiyah-macdonald-commutative-algebra:page-0092}
For a maximal ideal $\m$ in a Noetherian ring, an ideal $\q$ is $\m$-primary iff
$\rad(\q)=\m$ iff $\m^n\subseteq\q\subseteq\m$ for some $n$.
\uses{am-prop-4-2,am-prop-7-14}
\end{corollary}
\begin{proof}
The first equivalence is the definition and Proposition~\ref{am-prop-4-2};
the second follows by taking radicals and Proposition~\ref{am-prop-7-14}.
\end{proof}
\begin{proposition}[Associated primes in the Noetherian case]
\label{am-prop-7-17}
\dcref{7.17}
\source{atiyah-macdonald-commutative-algebra:page-0093}
For a proper ideal $\aideal$ in a Noetherian ring, the primes belonging to
$\aideal$ are exactly the radicals of $(\aideal:x)$ as $x$ ranges over $A$.
\end{proposition}
\begin{proof}
Pass to $A/\aideal$ and a minimal primary decomposition of zero.  For each
component, the chain of intersections of the other components eventually
annihilates; a last nonzero element has annihilator equal to the component's
radical.  Conversely Theorem~\ref{am-thm-4-5} identifies every radical of an
annihilator with an associated prime.
\uses{am-thm-4-5,am-prop-7-14}
\end{proof}
